% SHARD-ID: LB-06-MEMORY-QUANTIZATION
% SHARD-TITLE: Memory: what is refuted, what is exact, and what quantizes
% SHARD-SOURCES: theory/memory-quantization.md; theory/memory-quantization-general.md;
%   theory/mq-e.md; theory/corner-b-draft.md sections 2, 5, 6, 7, 10;
%   claims/CLAIMS.md rows M, M-flux, B3, M-quant, M-quant-G, Mq-E, Mq-AD3,
%   M-tk, N2, Bc; numerics/results/memory-scan-1.json,
%   numerics/results/spin1-bc-falsifier.json.
% SHARD-CLAIMS: M, M-flux, B3, M-quant, M-quant-G, Mq-E, Mq-AD3, M-tk, N2, Bc.

% Shard-local macro: the unit shift on l^2(Z), to keep the translation
% operator of the Fano Jacobi matrix visually distinct from the
% transmission probability T(k) of Definition (kink-magnon scattering data).
\newcommand{\shiftop}{\mathsf{T}}

\section{Memory: what is refuted, what is exact, and what quantizes}
\label{sec:memory-quantization}

The memory corner began with an attractive guess: that the wall displacement
should equal the zero-frequency limit of the soft factor, a lattice analogue of
the Braginsky--Thorne relation. That guess is false, and saying exactly why it
is false is the most useful thing this section does. What survives is sharper
and, in a sense, better: an exact flux identity that holds at finite time and
finite volume with no hypotheses at all, a charge ledger that constrains the
displacement channel by channel, and --- conditional on the scattering
hypothesis \cref{def:h-ad} --- a genuinely quantized displacement whose quantum
is a charge, not a phase. Soft data re-enter at exactly one place: the value of
the transmission probability $T(k)$.

\subsection{The refuted conjecture}

\begin{refutedclaim}[The soft-factor reading of the memory]
\label{ref:soft-factor-memory}
\statusRefuted\ \ \emph{Withdrawn statement.} The wall displacement $\delta x$
equals the zero-frequency limit of the soft factor $\softS$ summed over the
event --- the lattice Braginsky--Thorne relation as written in the campaign
brief.

\smallskip
\emph{Why it fails.} Two independent mechanisms kill it in the model of
\cref{def:kink-model}.
\begin{enumerate}
\item Once the channels of \cref{def:h-ad} exist, conditional charge
  bookkeeping fixes the \emph{same} per-channel displacement for every $k$,
  every $\Delta$ and every packet shape (\cref{thm:M-quant}), whereas any
  soft-factor expression varies with all three.
\item $\delta x$ is insensitive to the transmission phase $\delta_t(k)$
  entirely. A purely transmitting wall with $\delta_t\equiv0$ still displaces by
  exactly $-1/s$. The soft factor is a phase; the memory quantum is a charge,
  and the two are not the same kind of object.
\end{enumerate}

\smallskip
\emph{What survives.} Two statements, and no more. The flux identity
\cref{thm:M-flux}, which says that the memory is the DC weight of the
\emph{physical boundary current}; and the conditional quantization theorem
\cref{thm:M-quant}, into which soft data enter only through the transmission
probability $T(k)$. No reading in terms of virtual or bond data is claimed.
\end{refutedclaim}

\provenance{M}%
{disproved in \texttt{theory/memory-quantization.md} \S1 and
 \texttt{theory/corner-b-draft.md} \S\S2, 6, 10}%
{\texttt{theory/checks/mquant\_check.py} tests the flux identity and the
 empirical scan only}%
{\texttt{theory/verdicts/corpus-r2.md}}

\subsection{The exact flux identity}

\begin{theorem}[The memory is the DC weight of the physical boundary current]
\label{thm:M-flux}
\statusProved\ \ Let the chain be finite and open, let $\varrho(t)$ be any
state, let $W=[a,b]$ be any interior window, and assume only the local
continuity equation
\begin{equation}
  \frac{d}{dt}S^z_x=j_{x-1|x}-j_{x|x+1}.
  \label{eq:continuity}
\end{equation}
Then for every finite $t_i<t_f$,
\begin{equation}
  \varrho_{t_f}(\wallX_W)-\varrho_{t_i}(\wallX_W)
  =\frac{1}{2s}\int_{t_i}^{t_f}\!dt\,
   \Bigl[\varrho_t\bigl(j_{a-1|a}\bigr)-\varrho_t\bigl(j_{b|b+1}\bigr)\Bigr]
  =\frac{1}{2s}\Bigl[\tilde\jmath_{a-1|a}(0)-\tilde\jmath_{b|b+1}(0)\Bigr],
  \label{eq:flux-identity}
\end{equation}
where $\tilde\jmath(0)$ denotes the finite-time Fourier transform of the
corresponding bond current evaluated at zero frequency. No kink, spectral,
large-time or asymptotic-completeness hypothesis is used.
\end{theorem}

\noindent\emph{Idea of proof.} The windowed wall observable \eqref{eq:wallX} is
a bounded local observable at every finite volume, so its Heisenberg derivative
may be computed term by term. Inserting \eqref{eq:continuity} into the finite
sum, the interior currents cancel in pairs and only the two boundary bonds
survive:
$\dot{\wallX}_W=\tfrac{1}{2s}\bigl(j_{a-1|a}-j_{b|b+1}\bigr)$. Integrating from
$t_i$ to $t_f$ and reading the result as a Fourier transform at zero frequency
gives \eqref{eq:flux-identity}. The whole content is that $\wallX_W$ can change
\emph{only} through its two edges. \hfill$\square$

\begin{remark}[The current here is the physical one]
\label{rem:physical-current}
Nothing in \eqref{eq:flux-identity} identifies the boundary current with the
virtual bond potential $\mathcal J_b$ of the Noether analysis
(\cref{sec:symmetry-noether}). That identity writes the \emph{vacuum charge
density} as a difference of virtual insertions, in a different register. The
earlier claim that \eqref{eq:flux-identity} is a difference of virtual or bond
data is withdrawn; extending the virtual-potential reading to a kink-sector
dynamical theorem remains future work.
\end{remark}

\provenance{M-flux}%
{\texttt{theory/memory-quantization.md} \S1, Lemma Mq-flux, steps
 (1.1)--(1.4); \texttt{theory/corner-b-draft.md} \S2 Eq.\ (2.2)}%
{\texttt{theory/checks/mquant\_check.py}, named computation Mq-check F1: zero
 operator residue in \eqref{eq:continuity} and $3.37\times10^{-16}$
 finite-time residue in \eqref{eq:flux-identity} on frozen-boundary chains at
 $N=7,8$}%
{\texttt{theory/verdicts/corpus-r2.md}}

\subsection{Finite-time label rigidity and the charge ledger}

\begin{theorem}[The vacuum-pair label is rigid at finite time, and the charge
ledger holds]
\label{thm:B3}
\statusProved\ \ Assume a translation-invariant finite-range dynamics whose two
asymptotic vacua are stationary, and a state initially in the kink sector
$\Kk_{\alpha\beta}$. Then:
\begin{enumerate}
\item \emph{(Label rigidity.)} $\varrho_t\in\Kk_{\alpha\beta}$ for every finite
  $t$: finite-time dynamics fixes the vacuum-pair label $(\alpha,\beta)$, and
  the argument fixes the labels rather than merely preserving disjointness.
\item \emph{(Charge ledger.)} For an event separated in the sense of
  \cref{def:h-ad}, with an explicit cut $c\in W$ and leg-subtracted channel
  charges,
  \begin{equation}
    2s\cdot\delta x+\bigl(q_{\mathrm{out}}-q_{\mathrm{in}}\bigr)=0 .
    \label{eq:ledger}
  \end{equation}
\end{enumerate}
\end{theorem}

\noindent\emph{Idea of proof.} For rigidity, let $\alpha_t$ be the finite-range
Heisenberg automorphism. Lieb--Robinson quasi-locality supplies, for each local
pair $D,O$, each fixed $t$ and each $\varepsilon>0$, \emph{local} approximants
$D_\varepsilon,O_\varepsilon$ --- independent of the translation index $n$ ---
with $\norm{\alpha_t(D)-D_\varepsilon}<\varepsilon$ and
$\norm{\alpha_t(O)-O_\varepsilon}<\varepsilon$. Translation covariance gives
$\alpha_t(\tau_n(O))=\tau_n(\alpha_t(O))$, and stationarity of the vacua gives
$\omega_{\alpha/\beta}(\alpha_t(O))=\omega_{\alpha/\beta}(O)$. The two-factor
estimate
\begin{equation}
  \bigl|\varrho\bigl(\alpha_t(D)\alpha_t(\tau_n(O))\bigr)
  -\varrho\bigl(D_\varepsilon\tau_n(O_\varepsilon)\bigr)\bigr|
  \le\varepsilon\norm{O}+\bigl(\norm{D}+\varepsilon\bigr)\varepsilon
\end{equation}
then lets one take $n\to\mp\infty$ first, in the factorised local limit, and
$\varepsilon\downarrow0$ afterwards, reproducing the same left and right vacuum
functionals. For the ledger, a kink at coordinate $m$ contributes $2s(m-c)$ to
the regularised total charge, so the in and out totals are
$2s(m_i-c)+q_{\mathrm{in}}$ and $2s(m_f-c)+q_{\mathrm{out}}$; conservation of
the regularised charge and subtraction give \eqref{eq:ledger}. Specialising,
reflection has $q_{\mathrm{out}}=q_{\mathrm{in}}=-1$ and hence $\delta x=0$,
while transmission has $(q_{\mathrm{in}},q_{\mathrm{out}})=(-1,+1)$ and hence
$\delta x=-1/s$. \hfill$\square$

\begin{scope}[What this theorem does not say, and one thing it refutes]
\label{scope:B3}
The charges in \eqref{eq:ledger} are \emph{leg-subtracted} channel charges,
defined relative to the vacuum on which each separated leg sits. They are not
raw half-line charges. Indeed, the former raw half-line formula
$\delta x=(2s)^{-1}\Delta\mathfrak q^R_c=-(2s)^{-1}\Delta\mathfrak q^L_c$ is
\statusRefuted: in the transmitted spin-$1/2$ channel it predicts $+1$ (or
$-1$), while the ledger \eqref{eq:ledger} gives $-2$. Separately, the proposed
reading of the kink coordinate as a $\Z$-torsor, with a flat kink band, depends
on \cref{conj:K4} and is therefore a conjecture, not a consequence of this
theorem. The vacuum-pair label is rigid and is read off, not shifted; $\delta x$
records the leg-charge transfer within that fixed label.
\end{scope}

\provenance{B3}%
{\texttt{theory/corner-b-draft.md} \S7, steps (1.1)--(1.3), with the N9 density
 leaf supplied}%
{---}%
{\texttt{theory/verdicts/corpus-r2.md} (adjudicated with the N9 sweep)}

\subsection{Conditional quantization for the spin-\texorpdfstring{$1/2$}{1/2}
kink}

\begin{theorem}[Windowed memory is channel-quantized, conditional on the
scattering hypothesis]
\label{thm:M-quant}
\statusProvedConditional\ \ Let $\Delta>1$ and take the spin-$1/2$ Hamiltonian
of \cref{def:kink-model}, with unbroken $U(1)$ charge $S^z$ and asymptotic
densities $+1/2$ on the left and $-1/2$ on the right, selecting the kink sector
$\Kk_{\uparrow\downarrow}$. Prepare one right-moving magnon on the $\uparrow$
leg and use the windowed wall coordinate \eqref{eq:wallX}.

\textbf{Assume (H-AD), i.e.\ all four clauses AD1--AD4 of \cref{def:h-ad}, for
this scattering vector.} Then, in the limit order fixed by (AD4), the
asymptotic two-channel displacement observable obeys
\begin{equation}
  \Delta X=-\frac{1}{s}\,N_T,
  \qquad
  \delta x=-\frac{\braket{N_T}}{s},
  \label{eq:Mquant}
\end{equation}
and consequently
\begin{equation}
  \spec(\Delta X)\subseteq\Bigl\{-\frac1s,\,0\Bigr\},
  \qquad
  \delta x\in\Bigl[-\frac1s,0\Bigr],
  \qquad
  \operatorname{Var}_\Psi(\Delta X)
  =\frac{1}{s^2}\braket{N_T}\bigl(1-\braket{N_T}\bigr).
  \label{eq:Mquant-spec}
\end{equation}
The channel statement is independent of $k$, of $\Delta$, of the packet shape
and of the scattering phase. For the concrete spin-$1/2$ application $s=1/2$,
so $\delta x=-2\braket{N_T}$ and the variance is $4\braket{N_T}(1-\braket{N_T})$.
\end{theorem}

\noindent\emph{Idea of proof.} Choose a cut $c$ and regularise the total charge
by subtracting $+s$ to the left of $c$ and $-s$ to the right; a kink whose
windowed coordinate is $m$ then contributes $2s(m-c)$. The regularised charge
is exactly conserved, because the bond currents telescope and the frozen
endpoints carry no hopping current, and (AD1) with (AD4) carries that fixed
label into the in and out channel limits. On the reflected branch the in and out
totals are $2s(m_i-c)-1$ and $2s(m_R-c)-1$, so $m_R=m_i$; on the transmitted
branch they are $2s(m_i-c)-1$ and $2s(m_T-c)+1$, so $2s(m_T-m_i)=-2$ and
$m_T-m_i=-1/s$. Hence on the two outgoing channel spaces the displacement is
$0\cdot P_L-(1/s)P_T$, and conjugating by $W_+$ gives $\Delta X=-(1/s)N_T$ on
$\Hilb_{\mathrm{sc}}$. Since $N_T$ is an orthogonal projection, its spectrum
lies in $\{0,1\}$ and $N_T^2=N_T$, which gives \eqref{eq:Mquant-spec}
immediately. \hfill$\square$

\begin{scope}[Exactly what is conditional, and what the variance is a variance
of]
\label{scope:M-quant}
The hypothesis (H-AD) is used at precisely three places: to say that dressing
inside a fixed-charge kink sector does not change the kink eigenvalue, to
replace interacting late-time states by separated channels, and to package the
channel weights into the projection $N_T$. The flux identity
\cref{thm:M-flux} and the charge arithmetic use none of it. (H-AD) for the
unprojected chain remains open; it is proved only for the projected
three-wall component (\cref{cor:Mq-AD3}).

The abstract arithmetic behind \eqref{eq:Mquant} is only the ledger
\eqref{eq:ledger}; it does not by itself establish the required channels for a
higher-spin $XXZ$ chain, nor for a general triple (compact group, injective
MPS, finite-range Hamiltonian). That generality gap is closed only in the
conditional form of \cref{thm:M-quant-G}.

It is the \emph{outcomes} that are quantized, not every expectation value. The
expectation $\delta x=-\braket{N_T}/s$ is the convex combination selected by
$\braket{N_T}\in[0,1]$. Likewise $\operatorname{Var}(\Delta X)$ is the variance
of the two-channel displacement observable; identifying it with the one-time
variance $\operatorname{Var}(\wallX_W)$ would additionally require a sharp
initial wall eigenstate and an explicit two-time measurement convention, and
neither is implicit here.
\end{scope}

\begin{remark}[Finite-window error budget]
\label{rem:error-budget}
At finite window and finite separation it is useful to keep the theorem and the
measurement apart:
\begin{equation}
  \Bigl|\delta x_W+\frac{\braket{N_T}}{s}\Bigr|
  \le\epsilon_{\mathrm{AD}}
   +C_{\tilde\lambda}\tilde\lambda^{d_W}
   +\frac{N_W(t_i)+N_W(t_f)}{s}
   +\epsilon_{\mathrm{num}},
  \label{eq:budget}
\end{equation}
with $\tilde\lambda\in(\lambda_E,1)$, $d_W$ the minimum kink-core-to-edge
distance, $\epsilon_{\mathrm{AD}}$ an empirical local charge-decomposition
defect incurred when the exact channels are approximated numerically, $N_W$ the
residual leg content inside $W$, and $\epsilon_{\mathrm{num}}$ the estimator and
evolution errors. All terms vanish in the exact conditional limits, and
$\epsilon_{\mathrm{AD}}$ is not a return to the superseded norm-mixture
formulation of \cref{rem:h-ad-r1}. A projection error of order $\Delta^{-2}$ is
added only when the full $\braket{N_T}$ is replaced by the Fano value of
\cref{thm:M-tk}; it is not an error in \eqref{eq:Mquant}.
\end{remark}

\provenance{M-quant}%
{\texttt{theory/memory-quantization.md} \S4 steps (1.1)--(1.7) and \S5
 Corollary Mq-quant; \texttt{theory/corner-b-draft.md} \S6}%
{\texttt{theory/checks/mquant\_check.py} tests the flux identity and the
 empirical scan only; (H-AD) is an assumption and is tested by no gate}%
{\texttt{theory/verdicts/corpus-r2.md} (proved conditional on
 \cref{def:h-ad}); full-chain (H-AD) tracked as open work}

\subsection{The compact-group generalisation}

\begin{theorem}[The memory law for compact $G$ and injective MPS vacua]
\label{thm:M-quant-G}
\statusProvedConditional\ \ The following is proved \emph{only} as a
conditional implication: the hypothesis package H-MQG(1)--(4) together with
H-AD-G implies the conclusions \eqref{eq:G1}--\eqref{eq:G2}.

\smallskip
\noindent\textbf{Hypotheses H-MQG.}
\begin{enumerate}
\item[(1)] $G$ is compact and $\{A_\gamma\}_{\gamma\in\Omega_{\mathrm{vac}}}$ is
  a $G$-covariant family of injective canonical MPS tensors of one common,
  arbitrary finite bond dimension $\chi$ (\cref{sec:mps-setting}). Fix $\alpha$
  and $\beta=g\cdot\alpha\neq\alpha$, and use the kink sector
  $\Kk_{\alpha\beta}$ supplied by \cref{sec:corner-a}.
\item[(2)] A common unbroken compact abelian subgroup lies in
  $H_\alpha\cap H_\beta$. Choose one primitive circle direction $\xi$ in its Lie
  algebra and fix, once for both tails, the Hermitian on-site charge
  $S^z:=-i\,q(\xi)$. In this fixed convention $\omega_\alpha(S^z)=+s$ and
  $\omega_\beta(S^z)=-s$, with $s>0$. For a higher-dimensional torus the
  theorem is applied componentwise after choosing $\xi$; a finite abelian group
  with no circle direction is not enough, since it supplies no local current.
\item[(3)] $H$ is translation invariant, finite range and $G$-invariant in the
  exact sense of \cref{sec:symmetry-noether}, and the two vacua are stationary.
  Its infinite-volume dynamics therefore obeys the cut-current continuity
  equation for the selected charge.
\item[(4)] The selected incoming state is a \emph{fixed} wave packet in the
  $\ell^1$ class $\Kk^{(1)}_{\alpha\beta}$ of \cref{def:l1-kink-class}, and its
  wall coordinate is exactly the windowed observable \eqref{eq:wallX} with the
  charge and $s$ of item (2). No plane wave is admitted in place of this packet.
\end{enumerate}

\smallskip
\noindent\textbf{Hypothesis H-AD-G.} This is shorthand, not a replacement
definition, for all of \cref{def:h-ad}(AD1)--(AD4) holding for that vector, with
one incoming left leg, one reflected outgoing leg and one transmitted outgoing
leg, whose charges --- measured relative to the vacuum supporting each leg, in
the single convention of item (2) --- are
\begin{equation}
  q_{\mathrm{in}}=-1,\qquad q_L=-1\ \text{(reflected)},\qquad q_T=+1 .
\end{equation}
(AD2)'s transmitted-channel projection is $N_T$; (AD1) excludes any further
propagating channel, and the selected vector has no bound-state component.

\smallskip
\noindent\textbf{Conclusion.} Form the infinite-volume dynamics and the wave
operators first; fix the packet; at each fixed window take the incoming and
outgoing large-time limits; only then increase $W=[a,b]$ to $\Z$, with both
kink-core distances tending to infinity. In this order,
\begin{align}
  \Delta X&=-\frac{1}{s}N_T,
  &\delta x&=-\frac{\braket{N_T}}{s},
  \label{eq:G1}\\
  \spec(\Delta X)&\subseteq\Bigl\{0,-\frac1s\Bigr\},
  &\operatorname{Var}(\Delta X)
   &=\frac{1}{s^2}\braket{N_T}\bigl(1-\braket{N_T}\bigr).
  \label{eq:G2}
\end{align}
The outcome quantum is fixed by charge conservation alone: reflection transfers
no charge between the two vacuum conventions, while transmission changes the
separated leg charge by $q_T-q_{\mathrm{in}}=2$. More generally, if that
difference is some definite $\nu$, the same proof gives the quantum
$-\nu/(2s)$; the displayed $-1/s$ therefore includes the primitive channel
normalisation of H-AD-G, and is not invariant under relabelling a charge-$\nu$
particle as charge one.
\end{theorem}

\noindent\emph{Idea of proof.} Four steps. First, the selected state has a
fixed asymptotic label $(\alpha,\beta)$ which the dynamics does not move, by
the rigidity clause of \cref{thm:B3}; so $\delta x$ is motion \emph{within} a
fixed sector, not a change of vacuum-pair label. Second, an exact finite-window
identity: for a cut $c\in W$,
\begin{equation}
  \sum_{x=a}^{c}\bigl(S^z_x-s\bigr)+\sum_{x=c+1}^{b}\bigl(S^z_x+s\bigr)
  =2s\bigl(\wallX_W-c\bigr),
  \label{eq:G3}
\end{equation}
which is pure finite algebra, valid for arbitrary states and arbitrary $\chi$,
since both sides expand to $\sum_{x=a}^{b}S^z_x+s(a+b-1-2c)$. Third, the two
tail series converge absolutely on the $\ell^1$ class, the Heisenberg
derivative of the window charge is the difference of the two physical cut
currents, and the ordered limits of (AD3)--(AD4) carry the conserved total ---
kink plus separated leg --- into the channels. Fourth, equating in and out
totals gives the ledger \eqref{eq:ledger}, which returns $\delta x_R=0$ and
$\delta x_T=-1/s$; packaging these by (AD1)--(AD2) gives
$\Delta X=-(1/s)W_+P_TW_+^*$ and hence \eqref{eq:G1}--\eqref{eq:G2}, without
replacing the coherent outgoing state by a mixture. \hfill$\square$

\begin{scope}[What the generalisation does and does not buy]
\label{scope:M-quant-G}
The theorem is proved \emph{only} as the implication stated. H-AD-G itself is
not proved; neither is the sector reduction \cref{thm:Mq-E} in any general
setting, nor (H-AD) for the unprojected chain. No soft zero is proved here: if
a separate result shows $\braket{N_T}\to0$, then \eqref{eq:G1} transfers that
zero to the memory, but nothing in this theorem produces it. No uniformity is
claimed --- not over packets, models, soft limits, a continuous vacuum
manifold, $G$, $\alpha,\beta$, $\chi$, tensors, Hamiltonians, times, core
decorations or deformations. At the general $\ell^1$ level the tail error is
the state-specific omitted sum; it tends to zero as $W\uparrow\Z$ with no rate.
Only if the asymptotic kink core is one fixed two-sided MPS decoration with
modifications in a fixed bounded support does one recover an exponential tail
$C_{\tilde\lambda}\tilde\lambda^{d_W}$, on exactly that fixed data.

The transitivity hypothesis (T) is \emph{not} required for
\eqref{eq:G1}--\eqref{eq:G2}. If it is imposed, it supplies the double-coset
label $\mathfrak d(\alpha,\beta)\in H_\alpha\backslash G/H_\alpha$ of the fixed
vacuum pair, and the scattering does not change that label. Finally, the
measured displacement is charge bookkeeping inside a fixed sector; it is not a
moment map, a flat kink torsor, an SPT index, or a deformation-independent
absolute charge.

Two convention warnings. Rescaling the generator rescales both $s$ and all leg
charges, so only $-(q_{\mathrm{out}}-q_{\mathrm{in}})/(2s)$ is convention
covariant. And a finite abelian symmetry alone is insufficient: without a
circle generator there is only a modular charge label, not a local continuous
current and not a real-valued displacement law.
\end{scope}

\begin{remark}[Consistency with the spin-$1/2$ case]
\label{rem:G-instantiation}
Taking $G=U(1)\rtimes\Z_2$, $\alpha=\uparrow$, $\beta=\downarrow$ and $s=1/2$
in \cref{def:kink-model}, the two tensors are injective product tensors with
$\chi=1$, so H-MQG(1)--(4) holds as the special case in which every
transfer-tail error vanishes exactly. With $q_{\mathrm{in}}=q_L=-1$ and
$q_T=+1$, \eqref{eq:G1}--\eqref{eq:G2} reduce to $\Delta X=-2N_T$,
$\delta x=-2\braket{N_T}$, $\spec(\Delta X)\subseteq\{0,-2\}$ and
$\operatorname{Var}(\Delta X)=4\braket{N_T}(1-\braket{N_T})$ --- exactly the
numbers of \cref{thm:M-quant}. A rank-two example needs a nonabelian ambient
group: a bare abelian group, or a selected central factor, would force
$\omega_\beta=\omega_\alpha$ on the selected charge and could not produce
opposite vacuum densities.
\end{remark}

\provenance{M-quant-G}%
{\texttt{theory/memory-quantization-general.md} \S2, steps (1.1)--(1.5), with
 the hypothesis package in \S0 and the scope discipline in \S4}%
{\texttt{theory/checks/mquant\_general\_check.py} (symbolic charge arithmetic
 and representation theory only; the red mode mutates $q_T$). It constructs no
 MPS tensors, no wave operators and no proof of H-AD-G}%
{\texttt{theory/verdicts/mquant-g-r2.md} (promoted as the conditional
 implication only)}

\subsection{The exactly solvable sector: the Fano graph}

\begin{theorem}[The projected three-wall component is a Fano graph]
\label{thm:Mq-E}
\statusProved\ \ Take the spin-$1/2$ model of \cref{def:kink-model} with $J>0$,
$\Delta>1$, kink orientation $\uparrow$ on the left and $\downarrow$ on the
right, and a fixed integer regularised charge label $\mu$. Let $P_3$ project
the fixed-$S^z$ kink-plus-one-magnon space onto configurations with one or
three domain walls, retain the connected component containing the three
consecutive walls $(\mu-1,\mu,\mu+1)$, and put $H_3=P_3H_{XXZ}P_3$. Then, for
every admissible frozen-boundary interval and at infinite volume, there is an
explicit unitary $U$ from that component onto
$\ell^2(\Z)\oplus\C\ket{d}$ with
\begin{equation}
  UH_3U^*=H_0+V,\qquad
  H_0=\Bigl[E_c-\frac{J}{2}\bigl(\shiftop+\shiftop^*\bigr)\Bigr]\oplus E_d,
  \qquad
  V=-\frac{J}{2}\bigl(\ket{0}\bra{d}+\ket{d}\bra{0}\bigr),
  \label{eq:fano}
\end{equation}
where $\shiftop$ is the unit shift on $\ell^2(\Z)$,
$E_c=\tfrac{3}{2}J\Delta$ is the uniform channel on-site energy, and
$E_d=E_c-J\Delta=\tfrac12 J\Delta$ is the energy of the side vertex. The
channel structure is explicit:
\begin{itemize}
\item the negative tail is the incoming and reflected leg, of charge $-1$,
  carrying a sharp kink at bond $m:=\mu+1$;
\item the positive tail is the transmitted leg, of charge $+1$, carrying a
  sharp kink at bond $m-2=\mu-1$;
\item the side vertex $\ket{d}$ is the pure sharp kink at bond $\mu$.
\end{itemize}
Moreover every fixed-window observable maps to two exact channel-kink constants
plus a finite-support graph remainder: for $O\in\alg_W$,
\begin{equation}
  U P_C O P_C U^*
  =o_L(O)\,\Pi^-_W+o_R(O)\,\Pi^+_W+F_W,
  \label{eq:local-obs-map}
\end{equation}
with $o_L(O)=\bra{K_{\mu+1}}O\ket{K_{\mu+1}}$,
$o_R(O)=\bra{K_{\mu-1}}O\ket{K_{\mu-1}}$, $\Pi^\mp_W$ the projections onto
sufficiently remote negative and positive tails, and $F_W$ supported on a finite
graph set. For the wall observable itself, $o_L(\wallX_W)=\mu+1$,
$o_R(\wallX_W)=\mu-1$ and $\bra{K_\mu}\wallX_W\ket{K_\mu}=\mu$: zero reflected
displacement, and transmitted displacement exactly $-2$.
\end{theorem}

\noindent\emph{Idea of proof.} In the $S^z$ product basis, a configuration is
labelled by its ordered wall bonds. In the fixed charge sector the unique
one-wall word has its wall at bond $\mu$, and every three-wall word is uniquely
$(a,a+\ell,\mu+\ell)$ with $a\le\mu-1$ and $\ell\ge1$ --- a rectangle of
labels. The exchange term toggles the two wall indicators neighbouring the bond
it acts on, so a three-wall vertex has a surviving edge only when two of its
walls are adjacent; in rectangle coordinates that is exactly the row $\ell=1$
and the column $a=\mu-1$. Those two rays form a single two-sided path meeting at
the junction $(\mu-1,\mu,\mu+1)$, the one-wall word $K_\mu$ is attached only to
that junction, and every other rectangle vertex is isolated. Reading the row as
negative labels and the column as positive labels gives the unitary, and the
Ising energies ($3J\Delta/2$ per three-wall word, $J\Delta/2$ for the single
wall) with uniform hopping $-J/2$ give \eqref{eq:fano}. The finite maps are
literally restrictions of the infinite one along every cofinal sequence, since
each word is a finite deviation of a sharp kink. Identity
\eqref{eq:local-obs-map} follows because a $W$-local operator has zero matrix
element between product words differing outside $W$, so all but finitely many
tail rows are diagonal with the two kink constants. \hfill$\square$

\begin{scope}[Projected, not full chain]
\label{scope:Mq-E}
This is a theorem about the compressed operator $H_3=P_3H_{XXZ}P_3$, not about
the unprojected $XXZ$ Hamiltonian. It does not remove the separate leakage
problem $P_3H_{XXZ}(1-P_3)\neq0$: five-wall and higher configurations are real,
not absent. Their weight is \emph{measured}, not bounded: full-sector exact
diagonalisation at $N=22$, $k=\pi/2$, packet width $\sigma=2.6$, over
$0\le Jt\le13$, gives $P(\ge5\ \text{walls})\approx8\times10^{-3}$ at
$\Delta=8$, $3\times10^{-2}$ at $\Delta=4$ and $1\times10^{-1}$ at $\Delta=2$,
consistent with $O(\Delta^{-2})$. That is evidence for a controlled lift, not a
uniform-in-time scattering estimate.
\end{scope}

\provenance{Mq-E}%
{\texttt{theory/mq-e.md} \S\S1--7 (enumeration, component, finite unitary,
 infinite-volume compatibility, labels and charges, local-observable map)}%
{\texttt{theory/checks/mq\_e\_check.py} (finite-volume falsifier, safe under
 \texttt{python3 -O}; the red mode shifts the right wall pair and exits 1).
 Retained enumeration check: $N=14$, wall at bond $6$, component of size $12$
 with degree histogram $\{1{:}3,\,2{:}8,\,3{:}1\}$}%
{\texttt{theory/verdicts/blitz-mq-e-r1.md} (0 fatal, 0 major)}

\begin{corollary}[The projected component satisfies the scattering hypothesis]
\label{cor:Mq-AD3}
\statusProved\ \emph{(for the projected component)}\ \ The projected incoming
at-most-three-wall component of \cref{def:kink-model} satisfies all four clauses
AD1--AD4 of \cref{def:h-ad}, for every incoming packet with smooth momentum
amplitude compactly supported in $k\in(0,\pi)$. This does not prove
\cref{def:h-ad} for the unprojected full chain.
\end{corollary}

\noindent\emph{Idea of proof.} By \cref{thm:Mq-E} the compressed operator is
\eqref{eq:fano}, whose perturbation $V$ has rank at most two and is therefore
trace class. Kato--Rosenblum then gives the wave operators $W_\pm$ with
$\ran W_\pm=\Hilb_{\mathrm{ac}}(H_3)$, which is (AD1). Reflection about the
junction splits off an odd free channel; the even channel together with
$\ket{d}$ is a half-line Jacobi matrix differing from the constant-coefficient
one in finitely many entries. With
\begin{equation}
  m_0(z)=\bra{0}\Bigl(z-E_c+\tfrac{J}{2}(\shiftop+\shiftop^*)\Bigr)^{-1}\ket{0}
  =\bigl[(z-E_c)^2-J^2\bigr]^{-1/2},
\end{equation}
every cyclic resolvent entry is rational in $z$ and $m_0(z)$ with denominator
$z-E_d-(J^2/4)m_0(z)$. Inside the band $[E_c-J,E_c+J]$ the boundary value
$m_0(E+i0)$ has nonzero imaginary part, so that denominator has no real zero and
the spectral measure is absolutely continuous there; outside the band its zeros
are isolated and finite in number, and at the two thresholds the Jacobi
solutions are constant or linear (up to a sign alternation) and are not square
summable. Hence there is no singular continuous spectrum and
$\Hilb=\Hilb_b\oplus\ran W_\pm$ with $\dim\Hilb_b<\infty$. Finally, stationary
phase for $E(k)=E_c-J\cos k$ shows a smooth packet supported away from
$k=0,\pi$ leaves every finite set as $t\to\pm\infty$; the wave operators
transfer that local decay to $H_3$, and \cref{thm:Mq-E} labels the left and
right tails as the reflected and transmitted channels with charges $-1$ and
$+1$. Those are (AD2)--(AD3), and the strong limits impose (AD4).
\hfill$\square$

\provenance{Mq-AD3}%
{\texttt{theory/memory-quantization.md} \S3, steps (1.1)--(1.5);
 \texttt{theory/mq-e.md} \S6}%
{\texttt{theory/checks/mq\_e\_check.py} tests the sector reduction, not the
 spectral theorems}%
{\texttt{theory/verdicts/blitz-mq-e-r1.md}}

\subsection{The transmission coefficient and its quadratic soft zero}

\begin{theorem}[Fano transmission and the quadratic zero]
\label{thm:M-tk}
\statusProved\ \emph{(projected, unconditional)}\ \ On the projected incoming
at-most-three-wall component of \cref{def:kink-model}, the kink--magnon
scattering data of \cref{def:scattering-data} are
\begin{equation}
  t(k)=\Bigl[1+\frac{iJ^2}{4\,\omega(k)\,v(k)}\Bigr]^{-1},
  \qquad
  r(k)=-\frac{iJ^2}{4\,\omega(k)\,v(k)}\,t(k),
  \qquad
  T(k)=\Bigl[1+\Bigl(\frac{J^2}{4\,\omega(k)\,v(k)}\Bigr)^{2}\Bigr]^{-1},
  \label{eq:tk}
\end{equation}
with transmission phase
$\delta_t(k)=-\arctan\bigl[J^2/(4\,\omega(k)\,v(k))\bigr]$ and $\omega,v$ as in
\eqref{eq:dispersion}. In the soft limit $k\to0$ at fixed $\Delta$, where
$v\to0$ and $\omega\to J(\Delta-1)$,
\begin{equation}
  T(k)=16(\Delta-1)^2k^2+O(k^4),
  \qquad
  R(k)=1-16(\Delta-1)^2k^2+O(k^4).
  \label{eq:soft-zero}
\end{equation}
A soft magnon is thus totally reflected, with a quadratic zero in the
transmission. Combining with \eqref{eq:Mquant},
$\delta x(k)=-16(\Delta-1)^2k^2/s+O(k^4)$.
\end{theorem}

\noindent\emph{Idea of proof.} By \cref{thm:Mq-E} the effective one-body
problem is a uniform tight-binding chain with hopping $-J/2$ and a single
side-coupled level --- the pure kink --- attached at the junction with
amplitude $-J/2$ and detuning $\omega(k)$. Eliminating the side level leaves an
energy-dependent single-site potential $U(k)=(J/2)^2/\omega(k)$ at the
junction. Matching $\psi_n=e^{ikn}+re^{-ikn}$ for $n\le0$ to $\psi_n=te^{ikn}$
for $n\ge0$ across that impurity gives $-2i\tau r\sin k+Ut=0$ and $t=1+r$; with
$2\tau\sin k=J\sin k=v(k)$ this is \eqref{eq:tk}. \hfill$\square$

\begin{remark}[Two useful limits]
\label{rem:tk-limits}
For large $\Delta$ (the Ising regime), $J^2/(4\omega v)\approx1/(4\Delta\sin k)$,
so $R(k)\approx1/(16\Delta^2\sin^2k)$: the wall is transparent up to
$O(\Delta^{-2})$, the lattice counterpart of the reflectionless continuum domain
wall. In the frequency variable, using
$\omega-\omega_{\mathrm{gap}}=J(1-\cos k)\approx Jk^2/2$, the soft law
\eqref{eq:soft-zero} reads
$T\approx32(\Delta-1)^2(\omega-\omega_{\mathrm{gap}})/J$ --- linear in the
excess energy above the gap. The crossover momentum is $k_*=1/(4(\Delta-1))$.
\end{remark}

\begin{scope}[Projected and non-universal]
\label{scope:M-tk}
Equations \eqref{eq:tk}--\eqref{eq:soft-zero} are exact for the Fano graph, and
therefore unconditional on the projected component, because \cref{thm:Mq-E}
identifies that graph with the all-volume incoming component. They are
\emph{not} exact for the unprojected $H_{XXZ}$, because of the measured
$O(\Delta^{-2})$ leakage into five-wall and higher configurations. The
unprojected formula, and the universality of both the quadratic zero and its
coefficient $16(\Delta-1)^2$, remain \statusConjecture. What a genuine soft
theorem would have to supply is that this zero and its coefficient are functions
of the kink's asymptotic data --- vacuum pair, $U(1)$ charge, gap --- and not of
the microscopic tensor. That is the real obligation on the edge from the soft
corner to the memory corner, and it is not discharged here.

A separate warning, worth repeating: $d\delta_t/dk$ is the Wigner--Eisenbud
spatial shift of the \emph{transmitted magnon}. It is smooth and non-quantised
and it is not the wall displacement. The leakage affects $T(k)$; it does not
affect the memory law \eqref{eq:Mquant}, which is a conservation law.
\end{scope}

\provenance{M-tk}%
{\texttt{theory/corner-b-draft.md} \S5.2 Eqs.\ (5.1)--(5.3);
 \texttt{theory/memory-quantization.md} \S5 Eqs.\ (Mq.9)--(Mq.10);
 \texttt{theory/mq-e.md} \S7 step (1.2)}%
{\texttt{theory/checks/mq\_e\_check.py};
 \texttt{theory/checks/crosscheck\_corner\_b\_tk.py} remains an empirical
 full-chain cross-check only}%
{\texttt{theory/verdicts/blitz-mq-e-r1.md}}

\subsection{What the numerics see}

\begin{observation}[The wavepacket scan reproduces the conditional law]
\label{obs:N2}
\statusSketch\ \ An $XXZ$ wavepacket scan matches the spin-$1/2$ prediction
$\delta x=-2\braket{N_T}$ of \cref{thm:M-quant} at the recorded precision. On
the nine clean rows at $N=160$, three-wall truncation, standoff $36$, with
trapped weight below $10^{-6}$:
\begin{itemize}
\item the headline estimator satisfies $\max|\delta x_1+2T|=0.004330$ sites,
  i.e.\ $0.004$ at the reported precision;
\item the integrated-magnetisation estimator $\delta x_2$, which is the direct
  proxy for \eqref{eq:wallX}, satisfies $\max|\delta x_2+2T|=0.001233$ sites;
\item the largest estimator spread is $0.005563$ sites and the largest possible
  trapped-charge term is $2\times(8.862\times10^{-7})$ sites, so both residuals
  lie inside the scan's own budget \eqref{eq:budget};
\item norm and energy drifts are at most $7.75\times10^{-13}$ and
  $6.08\times10^{-10}$, subleading to the measurement geometry.
\end{itemize}
This is \statusSketch\ and stays there: it is numerical evidence for the
conditional theorem, not a theorem, and in particular not a proof of
\cref{def:h-ad} for the full chain. Comparing three-wall against five-wall
truncation shifts the robust integrated estimator by $0.0096$ sites at
$\Delta=2$, $N=56$, about half a percent, which together with the measured
$O(\Delta^{-2})$ higher-wall probability supports but does not establish the
full-chain lift.
\end{observation}

\provenance{N2}%
{no proof; empirical scan only, reported in
 \texttt{theory/memory-quantization.md} \S6, named computation Mq-check N1}%
{\texttt{numerics/results/memory-scan-1.json}: $\delta x\approx-2T$ within
 $0.004$ sites}%
{\texttt{theory/verdicts/corpus-r2.md} (held at sketch)}

\subsection{Are the two twos the same number?}

\begin{conjecture}[The soft slope and the memory quantum are one charge datum]
\label{conj:Bc}
\statusConjecture\ \ The soft phase slope of the two-magnon problem and the
memory channel quantum of \cref{thm:M-quant} are the same asymptotic-charge
datum, namely $|q_{\mathrm{hard}}|/s$, where $q_{\mathrm{hard}}$ is the hard
leg's $U(1)$ charge relative to its vacuum and $s$ is the vacuum density of
\cref{def:memory-observables}.
\end{conjecture}

\noindent\emph{Where the conjecture comes from.} In the isotropic ferromagnet
the two-magnon phase has $d\delta_{\mathrm{phys}}/dk_s|_0=2$ on the half-zone
(\cref{sec:two-magnon}), i.e.\ a soft Wigner shift of two sites off a hard
magnon; and the kink memory quantum of \cref{thm:M-quant} is $1/s=2$ sites at
$s=1/2$. Two independent-looking calculations return the same number. The
conjecture is that this is not a coincidence.

\smallskip
\noindent\emph{Evidence, stated as evidence.} A pre-registered falsifier was run
and \emph{survived}. Across site spins $s\in\{1/2,1,3/2,2\}$, the measured soft
slope is $1/s$ and the measured ratio is $\delta x/N_T=-1/s$, both inside
pre-registered $8\%$ criteria, on an ansatz-free ring plus dynamics computation
with a Bethe residual below $1.1\times10^{-14}$. A second, sharper falsifier
targeted the charge-two case: three-magnon exact diagonalisation at $N=100$
measured $4.07001\pm0.08981$ against the predicted $|q_{\mathrm{hard}}|/s=4$ at
$s=1/2$, inside the pre-registered absolute window of $0.35$, with the
dense-sector oracle green and the hopping mutation red.

\smallskip
\noindent\emph{Why this is not a proof.} Both measurements are finite-packet
evidence. There is no extrapolation in system size or packet width, and no
charge-two kink-memory test at all --- the memory half of the conjecture has
been checked only at $|q_{\mathrm{hard}}|=1$. The soft half is generalised to
every site spin by the exact two-body slope theorem of \cref{sec:two-magnon},
but the memory half rests on the conditional theorems \cref{thm:M-quant} and
\cref{thm:M-quant-G}, and the factor at $|q_{\mathrm{hard}}|>1$ is open. No
cross-reference between the two halves promotes the conjecture; it remains
\statusConjecture. Should a spin-$1$ two-magnon phase ever return a slope of $2$
rather than $1$, the coincidence is numerology and must be dropped.

\provenance{Bc}%
{no proof shard; statement recorded in \texttt{theory/corner-b-draft.md} \S10
 and \S5.3}%
{\texttt{theory/checks/spin\_s\_slope\_check.py};
 \texttt{numerics/results/spin1-bc-falsifier.json} (84 ring, 16 dynamics and
 22 memory rows); \texttt{numerics/results/spin1-bc-crosscheck.json};
 \texttt{numerics/test/test\_spins\_twomagnon.jl} and
 \texttt{test\_spins\_memory.jl} (decision tests with fixed criteria);
 \texttt{theory/lanes/blitz-2026-08-29/bc-charge2/RESULT.md} with
 \texttt{test\_bc\_charge2.jl} and its green, red-stub and red-mutation logs}%
{falsifier survived 2026-08-26 and again on the charge-two lane 2026-08-29;
 status unchanged}
